\documentclass[spanish,letterpaper,oneside]{article}
\def\documenttitle {Delitos contra la libertad personal}
\def\documentsubtitle {Libertad, custodia y clasificación jurídica}
\def\documentsubject {Actividad 5 - Derecho penal especial mexicano}
\def\documentauthor {Martín Jonathan de la Cruz}
\def\coursename {Derecho penal especial mexicano}
\def\coursecode {020}
\def\universityname {Universidad Abierta y a Distancia de México}
\def\universityfaculty {Licenciatura en Derecho}
\def\universitydepartment {Derecho penal especial mexicano}
\def\universitydepartmentimage {departamentos/UnADM}
\def\universitydepartmentimagecfg {height=1.57cm}
\def\universitylocation {Roma Norte, Ciudad de México}
\def\authortable {
  \begin{tabular}{ll}
    \textbf{Alumno:} & Martín Jonathan de la Cruz \\
    Matrícula: & ES2611202040 \\
    Figura docente: & Verónica Ortiz López \\
    Grupo: & 020 \\
    Semestre/Bloque: & 2 / 2 \\
    Semana: & 5 \\
    \multicolumn{2}{l}{Fecha: 23 de octubre de 2026} \\
    \multicolumn{2}{l}{\universitylocation}
  \end{tabular}
}
\input{template}
\def\indexstyle{}
\def\nameltcont{Tabla de contenido}
\setcounter{tocdepth}{2}
\usepackage{helvet}
\renewcommand{\familydefault}{\sfdefault}
\usepackage[most]{tcolorbox}
\usepackage{tabularx}
\usepackage{tikz}
\usetikzlibrary{babel}
\definecolor{casoverde}{HTML}{174A3A}
\definecolor{casopapel}{HTML}{F6F7F2}
\setcitestyle{authoryear,open={(},close={)},citesep={;}}
\def\pagemargintop{2.3}
\def\pagemarginbottom{2.3}
\def\predocpageromannumber{true}
\def\predocresetpagenumber{true}
\newcommand{\insertcoverwatermark}{%
  \AddToShipoutPictureBG*{%
    \begin{tikzpicture}[remember picture,overlay]
      \node[opacity=0.13,inner sep=0pt,yshift=-0.35cm]
        at (current page.center) {\includegraphics[width=0.72\paperwidth]{img/departamentos/UnADM.pdf}};
    \end{tikzpicture}%
  }%
}
\newtcolorbox{casoestudio}[1]{
  enhanced,colback=casopapel,colframe=casoverde,
  boxrule=0pt,leftrule=2.4pt,arc=1mm,
  left=2.4mm,right=2.4mm,top=1.7mm,bottom=1.7mm,
  title={#1},colbacktitle=casoverde,coltitle=white,
  fonttitle=\bfseries,fontupper=\fontsize{12}{13.8}\selectfont,
  before skip=5pt,after skip=5pt
}
\newcommand{\fichacaso}[1]{%
  \begin{tcolorbox}[colback=white,colframe=casoverde!25,
    boxrule=0.4pt,arc=0.6mm,left=1.6mm,right=1.6mm,top=1mm,bottom=1mm,
    before skip=0pt,after skip=3pt]
    #1
  \end{tcolorbox}%
}
\setlength{\parindent}{0pt}
\setlength{\parskip}{3pt}
\setlength{\tabcolsep}{4pt}
\renewcommand{\arraystretch}{1.08}

\newcommand{\tablaficha}[1]{%
  \begingroup\fontsize{10.4}{11.7}\selectfont
  \setlength{\tabcolsep}{3pt}\renewcommand{\arraystretch}{1.12}
  \begin{tabularx}{\linewidth}{@{}*{4}{>{\raggedright\arraybackslash}X}@{}}
  \toprule
  \textbf{Hechos conocidos} & \textbf{Elementos legales identificados} & \textbf{Datos que faltan} & \textbf{Posible figura y fundamento} \\
  \midrule
  #1
  \bottomrule
  \end{tabularx}
  \endgroup\par\vspace{4pt}
}
\begin{document}
\insertcoverwatermark
\templatePortrait
\templatePagecfg
\hypersetup{urlcolor=casoverde,citecolor=black,linkcolor=casoverde,linktoc=all}
\phantomsection\addcontentsline{toc}{section}{Resumen}
\begin{abstractd}
La privación de libertad, el secuestro y la desaparición comparten una restricción de la autonomía, pero exigen elementos distintos. La finalidad de rescate permite examinar el secuestro; la intervención estatal y el ocultamiento son decisivos para la desaparición forzada. En situaciones que involucran menores, el dinero, el parentesco o la falta de devolución no sustituyen la comprobación de edad, consentimiento, guarda y finalidad. La comparación de normas generales, federales y capitalinas permite formular hipótesis revisables para tres fichas de atención jurídica. La conclusión prioriza la búsqueda inmediata de la persona desaparecida y reserva la clasificación definitiva hasta esclarecer el vínculo estatal y las conductas posteriores a la privación.

\textbf{Palabras clave:} libertad personal; secuestro; desaparición; tráfico de menores; guarda y custodia.
\end{abstractd}
\renewcommand{\cftsecleader}{\cftdotfill{\cftdotsep}}
\renewcommand{\cftsubsecleader}{\cftdotfill{\cftdotsep}}
\templateIndex
\templateFinalcfg
\hypersetup{urlcolor=casoverde,citecolor=black,linkcolor=casoverde}
\fontsize{12}{13.8}\selectfont
\sectionfont{\color{casoverde}\fontsize{14}{16}\selectfont\bfseries}
\subsectionfont{\color{casoverde}\fontsize{11.3}{13}\selectfont\bfseries}
\titlespacing*{\section}{0pt}{7pt}{4pt}
\titlespacing*{\subsection}{0pt}{5pt}{3pt}
\section{Introducción}
\hypertarget{cuerpo-inicio}{}
Impedir salir, exigir dinero por liberar a alguien o desconocer su paradero compromete la libertad, pero no identifica el delito. La legalidad exige una ``ley exactamente aplicable al delito de que se trata'' (CPEUM, art.~14) \citep{cpeumS4Vigente}. Sin lugar ni fecha precisados, contrastar los códigos federal (CPF) y capitalino (CPDF) con las leyes generales de secuestro (LGS) y desaparición (LGD) exige verificar competencia y vigencia.

\section{Libertad personal, finalidad y vínculos de cuidado}
\subsection*{Elementos que distinguen las restricciones de libertad}
\phantomsection\addcontentsline{toc}{subsection}{Elementos que distinguen las restricciones de libertad}
La \textbf{libertad personal} protege el desplazamiento; la \textbf{tipicidad} exige todos los requisitos legales. La \textbf{clasificación preliminar} es una hipótesis revisable. La privación es \textbf{permanente} mientras continúa la restricción, sin exigir veinticuatro horas ni confundirla con delitos sucesivos \citep[pp.~67 y 409]{amuchateguiPenal2012}. El secuestro añade una \textbf{finalidad específica}, como rescate; la desaparición forzada, \textbf{intervención estatal} y ocultación o negativa posterior. En menores, \textbf{guarda y custodia}, patria potestad y \textbf{consentimiento} son diferentes: el parentesco no los acredita ni autoriza comerciar con la persona \citep{secuestroS5,desaparicionS5,cpdfS4Vigente}.

% Técnica principal: simulación #18; auxiliares: cuadro comparativo #1,
% estudio de casos #27 e informe #11. Tabla docente: figuras en filas,
% cuatro criterios en columnas. Se conserva orientación vertical para
% integrar el producto compuesto dentro de cuatro páginas de cuerpo.
\subsection*{Tabla de investigación: cinco figuras y sus diferencias}
\phantomsection\addcontentsline{toc}{subsection}{Tabla de investigación: cinco figuras y sus diferencias}
\begingroup\fontsize{10.3}{11.7}\selectfont
\renewcommand{\arraystretch}{1.12}
\begin{tabularx}{\linewidth}{@{}>{\raggedright\arraybackslash}p{2.5cm}>{\raggedright\arraybackslash}p{3.05cm}>{\raggedright\arraybackslash}X>{\raggedright\arraybackslash}X@{}}
\toprule
\textbf{Figura delictiva} & \textbf{Ley y artículo} & \textbf{Conducta descrita} & \textbf{Elemento que la diferencia} \\
\midrule
Privación de la libertad personal & CPF 364, I; CPDF 160. & Un particular impide a otra persona ejercer su libertad ambulatoria. & CPDF 160 exige ausencia de propósito de lucro, daño o perjuicio; varias horas bastan si se integra el tipo. \\
\midrule
Secuestro & LGS 9, I, a; 10, I, b y c. CPF 366 derogado. & Privar de libertad para obtener rescate o beneficio. & Finalidad adicional. Grupo de dos o más y violencia son agravantes si se acreditan. \\
\midrule
Desaparición forzada & LGD 27 y 28; alternativa distinta: 34. & Privación seguida de abstención o negativa a reconocerla o informar; el 28 regula ocultación o negativa informativa. & Servidor público o particular con autorización, apoyo o aquiescencia estatal. El 34 tipifica privación por particulares con finalidad de ocultar. \\
\midrule
Tráfico de menores & CPF 366 Ter y quáter; CPDF 169--170. & Traslado o entrega ilícitos; en CPDF 169, entrega para custodia definitiva. & CPF: menor de 16 años y condiciones del beneficio; incluye entrega nacional. CPDF: ilegalidad, consentimiento y finalidad de custodia. \\
\midrule
Retención y sustracción de menores o incapaces & CPDF 171--173. & Retener o apartar de guarda legítima; el 173 incluye ocultar por determinados familiares. & Relación con la víctima, falta de consentimiento y situación jurídica de patria potestad, tutela o custodia. \\
\bottomrule
\end{tabularx}
\par\vspace{2pt}\textit{Fuentes:} CPF, CPDF, LGS y LGD vigentes \citep{cpfS4Vigente,cpdfS4Vigente,secuestroS5,desaparicionS5}.
\endgroup

\clearpage
\subsection*{Fichas 1 y 2: finalidad y desaparición}
\phantomsection\addcontentsline{toc}{subsection}{Fichas 1 y 2: finalidad y desaparición}
\begin{casoestudio}{Ficha 1. Privación de libertad y secuestro}
\tablaficha{
Impedir la salida de una habitación durante varias horas. & Restricción efectiva por un particular; consentimiento y ausencia de justificación. & Medio de impedir la salida, duración, autoría y propósito. & Posible privación: CPF 364, I; CPDF 160, si no hay finalidad de lucro, daño o perjuicio. \\
\midrule
Varias personas trasladan por la fuerza a un comerciante y exigen dinero para liberarlo. & Privación y propósito de rescate; pluralidad y violencia. & Identidades, comunicaciones, enlace entre exigencia y cautiverio, forma de ejecución. & Posible secuestro: LGS 9, I, a; examinar agravantes del 10, I, b y c. \\
}
\textbf{Análisis.} En el primer supuesto, documentar cerraduras, amenazas, posibilidades reales de salir y testimonios; varias horas no vuelven atípica la restricción. En cambio, la exigencia ligada a la liberación orienta al \textbf{rescate}; no es necesario presumir que el dinero ya se pagó. Conservar mensajes, videos y declaración del comerciante para corroborar esa conexión. La jurisprudencia II.2o.P. J/9 P (11a.), registro 2028646, exige la finalidad específica para secuestro; proviene del Segundo Circuito y su obligatoriedad depende del ámbito jurisdiccional correspondiente \citep{scjnSecuestroS5}. La pluralidad y violencia no sustituyen ese elemento. No corresponde calcular penas sin fijar hechos, intervención y competencia (LGS 23) \citep{secuestroS5}.
\end{casoestudio}

\begin{casoestudio}{Ficha 2. Desaparición forzada}
\tablaficha{
La familia desconoce el paradero. Un testigo refiere que personas armadas se llevaron a su integrante; se ignora el vínculo estatal. & Privación relatada; para LGD 27 se requiere nexo estatal y posterior abstención o negativa a reconocer o informar. & Identidad de captores; autorización, apoyo o aquiescencia; actos de negativa u ocultamiento y secuencia. & Información insuficiente para confirmar desaparición forzada (LGD 27--28). Examinar art.~34 si se acredita finalidad de ocultar. \\
}
\textbf{Análisis.} Estar armado no demuestra pertenecer al Estado. Ampliar el testimonio sobre vehículos, distintivos, ruta y tiempo; contrastar cámaras, registros de detención y comunicaciones obtenidas legalmente. El desconocimiento familiar tampoco prueba por sí solo la negativa típica. La \textbf{desaparición cometida por particulares} exige finalidad de ocultar a la víctima o su suerte o paradero: no surge automáticamente de no probar el vínculo estatal (LGD 34). Sin embargo, la incertidumbre no suspende la búsqueda: la jurisprudencia 1a./J.~35/2021, registro 2023814, la exige con debida diligencia y presunción de vida \citep{desaparicionS5,scjnBusquedaS5}.
\end{casoestudio}

\clearpage
\subsection*{Ficha 3: entrega, retención y sustracción de menores}
\phantomsection\addcontentsline{toc}{subsection}{Ficha 3: entrega, retención y sustracción de menores}
\begin{casoestudio}{Ficha 3. Situaciones que involucran menores}
\tablaficha{
Una persona entrega a un menor a terceros a cambio de dinero; se ignoran las circunstancias. & Entrega y contraprestación; falta verificar ilegalidad, edad, calidad de sujetos y finalidad. & Edad exacta; quién consintió, sus facultades; motivo del pago, destino y carácter definitivo de la custodia. & Posible tráfico: CPF 366 Ter/quáter o CPDF 169, según sus requisitos y competencia; clasificación no confirmada. \\
\midrule
Un progenitor se lleva a su hija y no la devuelve a quien tiene legalmente su cuidado. & Posible sustracción o retención; parentesco y situación jurídica de cuidado. & Resoluciones vigentes, patria potestad, tutela, guarda, permiso inicial y obligación de devolver; lugar y tiempo. & Examinar CPDF 173 por parentesco; CPDF 171 corresponde a sujetos sin el vínculo previsto ni tutela. \\
}
\textbf{Análisis.} El \textbf{tráfico de menores} del CPF 366 Ter comprende menores de dieciséis años y también la entrega en territorio nacional; deben acreditarse la ilicitud y las condiciones del beneficio económico indebido. El 366 quáter contempla hipótesis de reducción, por lo que no basta trasladar una definición abreviada. En CPDF 169, la entrega ilegal debe destinarse a \textbf{custodia definitiva}; la falta de consentimiento agrava y la entrega sin beneficio tiene regulación propia. El dinero, aislado, no aclara esos extremos \citep{cpfS4Vigente,cpdfS4Vigente}.

Además, la \textbf{trata de personas} del artículo 10 de su ley general exige fines de explotación; no pueden inferirse de la sola entrega remunerada. También habría que esclarecer si existió una adopción ilegal: el artículo 27 regula esa conducta con sujetos y requisitos propios \citep{trataS5}. Respecto del progenitor, \textbf{retención} implica mantener a la hija fuera de la guarda debida; \textbf{sustracción}, apartarla de ella. La filiación no equivale a custodia, pero tampoco basta para atribuir el CPDF 173: revisar las condiciones que el artículo establece sobre patria potestad, tutela o resolución de guarda y custodia. No convertir un incumplimiento familiar indeterminado en delito sin verificar esos elementos \citep{cpdfS4Vigente}.
\end{casoestudio}

\subsection*{La finalidad y la protección orientan la clasificación}
\phantomsection\addcontentsline{toc}{subsection}{La finalidad y la protección orientan la clasificación}
Las cinco figuras comparten límites a la disponibilidad de la persona; sus diferencias impiden usar una misma prueba para todas. La \textbf{presunción de inocencia} preserva la exigencia de acreditar cada elemento, mientras el \textbf{interés superior de la niñez} orienta la protección inmediata de menores (CPEUM, arts.~4 y 20, B, I) \citep{cpeumS4Vigente}. Por tanto, distinguir:
\begin{itemize}[leftmargin=*,topsep=2pt,itemsep=2pt,parsep=0pt]
\item \textbf{Finalidad y ejecución:} rescate frente a simple restricción; ocultamiento frente a paradero desconocido. Las agravantes y atenuantes requieren sus propios hechos, no etiquetas.
\item \textbf{Vínculo y facultades:} aquiescencia estatal frente a actuación particular; consentimiento válido y guarda legítima frente a parentesco o pago. La autorización familiar no vuelve lícita cualquier entrega.
\end{itemize}
Así, distinguir hechos, elementos y datos faltantes permite revisar la hipótesis sin posponer la protección \citep{unadmPenalS5Material}.

\clearpage
\section{Conclusiones}
Considero que la ficha 2 permite investigar una privación de libertad, pero no confirmar desaparición forzada. La LGD, art.~27, exige vínculo estatal y posterior abstención o negativa a reconocerla o informar. Verificaría identidad de los captores y registros de detención. Sostengo que la búsqueda debe comenzar ``bajo la presunción de que está viva'', conforme a la jurisprudencia 1a./J.~35/2021; esperar una clasificación definitiva pospondría la protección urgente \citep{desaparicionS5,scjnBusquedaS5}.\footnote{Se utilizó ChatGPT como apoyo de inteligencia artificial para la organización del texto y la preparación de un borrador. El autor asume la revisión crítica y la responsabilidad académica de los análisis, el desarrollo y las conclusiones.}
\hypertarget{cuerpo-fin}{}

\clearpage
\begingroup\fontsize{10}{11.5}\selectfont
\setlength{\bibsep}{4pt}
\phantomsection
\begin{thebibliography}{10}
\bibitem[Amuchategui Requena(2012)]{amuchateguiPenal2012}
Amuchategui Requena, I.~G. (2012). \emph{Derecho penal} (4.ª ed.). Oxford University Press México.

\bibitem[Cámara de Diputados(2021)]{secuestroS5}
Cámara de Diputados del H. Congreso de la Unión. (2021). \emph{Ley General para Prevenir y Sancionar los Delitos en Materia de Secuestro, Reglamentaria de la Fracción XXI del Artículo 73 de la Constitución Política de los Estados Unidos Mexicanos} (reforma DOF 20 de mayo de 2021). \url{https://www.diputados.gob.mx/LeyesBiblio/pdf/LGPSDMS_200521.pdf}

\bibitem[Cámara de Diputados(2025a)]{desaparicionS5}
Cámara de Diputados del H. Congreso de la Unión. (2025a). \emph{Ley General en Materia de Desaparición Forzada de Personas, Desaparición Cometida por Particulares y del Sistema Nacional de Búsqueda de Personas} (reforma DOF 16 de julio de 2025). \url{https://www.diputados.gob.mx/LeyesBiblio/pdf/LGMDFP.pdf}

\bibitem[Cámara de Diputados(2025b)]{trataS5}
Cámara de Diputados del H. Congreso de la Unión. (2025b). \emph{Ley General para Prevenir, Sancionar y Erradicar los Delitos en Materia de Trata de Personas y para la Protección y Asistencia a las Víctimas de estos Delitos} (reforma DOF 14 de noviembre de 2025). \url{https://www.diputados.gob.mx/LeyesBiblio/pdf/LGPSEDMTP.pdf}

\bibitem[Cámara de Diputados(2026a)]{cpfS4Vigente}
Cámara de Diputados del H. Congreso de la Unión. (2026a). \emph{Código Penal Federal} (reforma DOF 13 de marzo de 2026). \url{https://www.diputados.gob.mx/LeyesBiblio/pdf/CPF.pdf}

\bibitem[Cámara de Diputados(2026b)]{cpeumS4Vigente}
Cámara de Diputados del H. Congreso de la Unión. (2026b). \emph{Constitución Política de los Estados Unidos Mexicanos} (reforma DOF 7 de octubre de 2026). \url{https://www.diputados.gob.mx/LeyesBiblio/pdf/CPEUM.pdf}

\bibitem[Congreso de la Ciudad de México(2026)]{cpdfS4Vigente}
Congreso de la Ciudad de México. (2026). \emph{Código Penal para el Distrito Federal} (reforma GOCDMX 26 de agosto de 2026). \url{https://www.congresocdmx.gob.mx/archivo-deae456e36ecb11eac420f0f1839cc70e9ea4685.pdf}

\bibitem[Segundo Tribunal Colegiado en Materia Penal del Segundo Circuito(2024)]{scjnSecuestroS5}
Segundo Tribunal Colegiado en Materia Penal del Segundo Circuito. (2024). \emph{Secuestro. No se configura este delito ante la ausencia del elemento subjetivo específico de finalidad que lo rige} [Jurisprudencia II.2o.P. J/9 P (11a.); registro 2028646]. \emph{Gaceta del Semanario Judicial de la Federación}, libro 36, tomo V, 4397. \url{https://sjf2.scjn.gob.mx/detalle/tesis/2028646}

\bibitem[Suprema Corte de Justicia de la Nación(2021)]{scjnBusquedaS5}
Suprema Corte de Justicia de la Nación. (2021). \emph{Desaparición forzada de personas. El derecho a no ser víctima de desaparición forzada comprende el derecho a la búsqueda como parte de su núcleo esencial} [Jurisprudencia 1a./J. 35/2021 (11a.), Primera Sala; registro 2023814]. \emph{Gaceta del Semanario Judicial de la Federación}, libro 7, tomo II, 1198. \url{https://sjf2.scjn.gob.mx/detalle/tesis/2023814}

\bibitem[UnADM(2026)]{unadmPenalS5Material}
Universidad Abierta y a Distancia de México. (2026). \emph{Material de apoyo S5: Delitos contra la libertad personal} [Material de curso]. Aula virtual. \url{https://aulavirtual.unadmexico.mx/mod/resource/view.php?id=217987}
\end{thebibliography}
\endgroup
\end{document}
